% ID: SEL-D0062-OB-E04
% ESTADO: cotejada. Segunda pasada visual y borrador compilado revisados.
% Fuente: Selectividad/sel_2010_matematicas_aplicadas/mat_apli/mat_apli_junio/mat_apli_junio_ra/EXAMEN-ANDALUCÍA-4.pdf
% SHA-256: 4e10e40d60aebdab7570fa77b9ce953d7bf6c9a4c8869a5d3f9b7ca3c907c154
% Página: 2; opción B; ejercicio original 4.

Se sabe que el tiempo de reacción a un determinado estímulo se distribuye según una ley Normal de media desconocida y desviación típica $0.2$ segundos.
\begin{enumerate}[label=\alph*)]
\item \textbf{(1.25 puntos)} Observada una muestra aleatoria de tamaño 25 se ha obtenido una media muestral de $0.3$ segundos. Obtenga un intervalo de confianza para la media de la población con un nivel de confianza del $94\%$.
\item \textbf{(1.25 puntos)} A un nivel de confianza del $90\%$, \textquestiondown{}cuál será el tamaño muestral mínimo si el error cometido es inferior a $0.05$?
\end{enumerate}
